% Part: first-order-logic
% Chapter: models-theories
% Section: theories

\documentclass[../../../include/open-logic-section]{subfiles}

\begin{document}

\olfileid{fol}{mat}{the}

\olsection{Voorbeelden van eerste-ordetheorieën}

\begin{ex}
De theorie van strikte lineaire ordeningen in de taal~$\Lang L_<$
wordt geaxiomatiseerd door de verzameling
\begin{align*}
\{\quad & \lforall[x][\lnot x < x], \\
& \lforall[x][\lforall[y][((x < y \lor y <
    x) \lor x = y)]], \\
& \lforall[x][\lforall[y][\lforall[z][((x < y
      \land y < z) \lif x < z)]]] \quad \}
\end{align*}
Dit axiomasysteem karakteriseert de beoogde !!{structure}s volledig:
iedere strikte lineaire ordening is er een model van. Omgekeerd geldt
dat als $R$ een lineaire ordening op een verzameling~$X$ is, de
!!{structure}~$\Struct M$ met $\Domain M = X$ en
$\Assign{<}{M} = R$ een model van deze theorie is.
\end{ex}

\begin{ex}
De theorie van groepen in de taal met $\Obj 1$ (een !!{constant}) en
$\cdot$ (een tweeplaatsige !!{function}) wordt geaxiomatiseerd door
\begin{align*}
& \lforall[x][\eq[(x \cdot \Obj 1)][x]]\\
& \lforall[x][\lforall[y][\lforall[z][\eq[(x \cdot (y \cdot z))][((x
          \cdot y) \cdot z)]]]]\\
& \lforall[x][\lexists[y][\eq[(x \cdot y)][\Obj 1]]]
\end{align*}
\end{ex}

\begin{ex}
De theorie van de Peano-rekenkunde wordt geaxiomatiseerd door de
volgende zinnen in de taal van de rekenkunde~$\Lang L_A$.
\begin{align*}
& \lforall[x][\lforall[y][(\eq[x'][y'] \lif \eq[x][y])]]\\
& \lforall[x][\eq/[\Obj 0][x']]\\
& \lforall[x][\eq[(x + \Obj 0)][x]]\\
& \lforall[x][\lforall[y][\eq[(x + y')][(x + y)']]]\\
& \lforall[x][\eq[(x \times \Obj 0)][\Obj 0]]\\
& \lforall[x][\lforall[y][\eq[(x \times y')][((x \times y) + x)]]]\\
& \lforall[x][\lforall[y][(x < y \liff \lexists[z][\eq[(z' + x)][y])]]]\\
\intertext{en alle zinnen van de vorm}
& (!A(\Obj 0) \land \lforall[x][(!A(x) \lif !A(x'))]) \lif \lforall[x][!A(x)]
\end{align*}
Omdat er oneindig veel zinnen van de laatste vorm zijn, is dit
axiomasysteem oneindig. Die laatste vorm heet het
\emph{inductieschema}. (Eigenlijk is het inductieschema iets
ingewikkelder dan we hier doen voorkomen.)

Het laatste axioma is een \emph{expliciete definitie} van~$<$.
\end{ex}

\begin{ex}
De theorie van zuivere verzamelingen speelt een belangrijke rol in de
grondslagen (en in de filosofie) van de wiskunde. Een verzameling is
zuiver als al haar !!{element}s eveneens zuivere verzamelingen zijn.
De lege verzameling geldt daarom als zuiver, maar een verzameling die
iets als !!a{element} heeft dat geen verzameling is, is niet zuiver.
De zuivere verzamelingen zijn dus precies die welke uitsluitend uit de
lege verzameling en zonder ``oerelementen'' worden opgebouwd, dat wil
zeggen zonder objecten die zelf geen verzamelingen zijn.

Het volgende kan worden beschouwd als een axiomasysteem voor een
theorie van zuivere verzamelingen:
\begin{align*}
& \lexists[x][\lnot \lexists[y][y \in x]]\\
& \lforall[x][\lforall[y][(\lforall[z](z \in x \liff z \in y) \lif 
\eq[x][y])]]\\
& \lforall[x][\lforall[y][\lexists[z][\lforall[u][(u \in z \liff
(\eq[u][x] \lor \eq[u][y]))]]]]\\
& \lforall[x][\lexists[y][\lforall[z][(z \in y \liff \lexists[u][(z \in
        u \land u \in x)])]]]\\
\intertext{en alle zinnen van de vorm} &
\lexists[x][\lforall[y][(y \in x \liff !A(y))]]
\end{align*}
Het eerste axioma zegt dat er een verzameling zonder !!{element}s
bestaat (dus dat $\emptyset$ bestaat); het tweede zegt dat voor
verzamelingen extensionaliteit geldt; het derde dat voor alle
verzamelingen $X$ en $Y$ de verzameling $\{X, Y\}$ bestaat; het vierde
dat voor iedere verzameling $X$ de verzameling $\cup X$ bestaat,
waarbij $\cup X$ de vereniging van alle elementen van $X$ is.

De laatstgenoemde !!{sentence}s heten gezamenlijk het
\emph{naïeve comprehensieschema}. Het zegt in wezen dat voor iedere
$!A(x)$ de verzameling $\Setabs{x}{!A(x)}$ bestaat---op het eerste
gezicht dus een waar, nuttig en misschien zelfs noodzakelijk axioma.
Het heet ``naïef'' omdat het deze theorie, naar blijkt, onvervulbaar
maakt: als men voor $!A(y)$ de formule $\lnot y \in y$ neemt, verkrijgt
men de !!{sentence}
\[
\lexists[x][\lforall[y][(y \in x \liff \lnot y \in y)]]
\]
en deze !!{sentence} wordt in geen enkele !!{structure} vervuld.
\end{ex}

\begin{ex}
Op het gebied van de \emph{mereologie} is de
\emph{deel-geheelrelatie} een fundamentele relatie. Net als bij
verzamelingen bestaan er theorieën van de deel-geheelrelatie die
verschillende (soms onderling strijdige) opvattingen van deze relatie
axiomatiseren.

De taal van de mereologie bevat één tweeplaatsig
predicaatsymbool~$\Obj P$, en $\Atom{\Obj P}{x, y}$ ``betekent'' dat
$x$ een deel van~$y$ is. Met deze interpretatie voor ogen heet
!!a{structure} voor deze taal een \emph{deel-geheelstructuur}.
Natuurlijk verdient niet iedere structuur voor één tweeplaatsig
predicaat werkelijk die naam. Om de ``deel-geheelrelatie'' te kunnen
weergeven, moet $\Assign{\Obj P}{M}$ aan bepaalde voorwaarden voldoen.
Die kunnen we als axioma's voor een theorie van de deel-geheelrelatie
vastleggen. Zo is de deel-geheelrelatie een partiële ordening op
objecten: ieder object is een deel (zij het een \emph{oneigenlijk}
deel) van zichzelf; geen twee verschillende objecten kunnen delen van
elkaar zijn; en een deel van een deel van een object is zelf een deel
van dat object. Merk op dat ``is een deel van'' in deze betekenis op
``is een deelverzameling van'' lijkt, maar niet op ``is een element
van'', want die laatste relatie is noch reflexief, noch transitief.
\begin{align*}
& \lforall[x][\Atom{\Obj P}{x,x}] \\
& \lforall[x][\lforall[y][((\Part{x}{y} \land \Part{y}{x})
      \lif \eq[x][y])]] \\
& \lforall[x][\lforall[y][\lforall[z][((\Part{x}{y} \land
        \Part{y}{z}) \lif \Part{x}{z})]]]\\
\intertext{Bovendien heeft ieder tweetal objecten een mereologische som (een object
  dat deze twee objecten als delen heeft en in dit opzicht minimaal is).}  &
\lforall[x][\lforall[y][\lexists[z][\lforall[u][(\Part{z}{u} \liff
        (\Part{x}{u} \land \Part{y}{u}))]]]]
\end{align*}
Dit zijn slechts enkele van de grondbeginselen van de
deel-geheelrelatie die metafysici bestuderen. Verdere beginselen worden
echter al snel moeilijk te formuleren of op te schrijven zonder eerst
enkele gedefinieerde relaties in te voeren. Zo beschouwen de meeste
metafysici die in mereologie geïnteresseerd zijn ook het volgende als
een geldig beginsel: wanneer een object~$x$ een strikt deel~$y$ heeft,
heeft het ook een deel~$z$ dat geen delen met~$y$ gemeen heeft, en wel
zo dat de fusie (mereologische som) van $y$ en $z$ gelijk is aan~$x$.
\end{ex}

\end{document}
